\documentclass[11pt,a4paper]{article}
\usepackage[utf8]{inputenc}
\usepackage[T1]{fontenc}
\usepackage[a4paper,margin=2.5cm]{geometry}
\usepackage{hyperref}
\usepackage{listings}

\setlength{\parindent}{0pt}
\setlength{\parskip}{0.6em}

\hypersetup{
  colorlinks=true,
  linkcolor=black,
  urlcolor=blue
}

\lstset{
  basicstyle=\ttfamily\small,
  breaklines=true,
  columns=fullflexible,
  frame=single,
  numbers=left,
  numberstyle=\tiny,
  xleftmargin=1em
}

\title{CSM Robot + Force\\Rapport de bilan}
\author{Harambat Martin\\Orhategaray Eneko\\Larregle Achille\\Kam Melvyn}
\date{18/03/2026}

\begin{document}

\begin{titlepage}
\centering
\vspace*{2cm}
{\Large Conception des syst\`emes m\'ecatroniques\par}
\vspace{1.2cm}
{\huge \textbf{CSM Robot + Force}\par}
\vspace{0.4cm}
{\Large Rapport de bilan du projet\par}
\vspace{2cm}

{\large Composition de l'\'equipe\par}
\vspace{0.5cm}
{\Large
Harambat Martin\par
Orhategaray Eneko\par
Larregle Achille\par
Kam Melvyn\par
}

\vfill

{\large Ann\'ee universitaire 2025--2026\par}
{\large Date de remise : 18/03/2026\par}
\end{titlepage}

\tableofcontents
\clearpage

\section{Introduction}
Le projet \emph{CSM Robot + Force} a pour objectif d'int\'egrer, au sein d'une m\^eme application C\#, un capteur de force exteroceptif et un bras robotique xArm6. Le besoin de la mati\`ere est double. Il s'agit d'abord de comprendre et de valider s\'epar\'ement les deux briques mat\'erielles, puis de proposer une application de synth\`ese capable d'associer une mesure de force \`a un mouvement robotique pilot\'e.

Dans le code disponible, cette int\'egration repose principalement sur le projet \texttt{RobotForceIntegration}, compl\'et\'e par le projet \texttt{RS232-PC} utilis\'e comme banc de test pour la liaison s\'erie. L'architecture r\'ealis\'ee est coh\'erente avec l'esprit du cahier des charges : la communication capteur passe par \texttt{SerialPort}, le pilotage du robot repose sur le r\'eseau IP et la DLL constructeur \texttt{xarm.dll}, et le couplage entre les deux est assur\'e par une logique temporelle dans l'IHM WinForms.

Le pr\'esent rapport a \'et\'e remani\'e pour correspondre plus directement \`a une remise acad\'emique. Il explicite ce qui est effectivement impl\'ement\'e dans le code, signale clairement les limites observables, et \'evite d'attribuer au projet des fonctionnalit\'es absentes, en particulier un PID complet en boucle ferm\'ee.

\section{Conformit\'e au cahier des charges}
Le cahier des charges demandait explicitement les \'el\'ements suivants dans le ZIP de rendu :
\begin{itemize}
  \item le PDF de bilan du travail, obtenu \`a partir de \texttt{Rapport\_CSM\_Robot\_Force.tex} via le script \texttt{build\_rapport\_pdf.bat} ;
  \item une description rapide de l'IHM, trait\'ee \`a la section \ref{sec:ihm} ;
  \item le couplage Robot/Force, trait\'e \`a la section \ref{sec:couplage} ;
  \item l'asservissement en force, discut\'e \`a la section \ref{sec:asservissement} en restant fid\`ele \`a l'\'etat r\'eel du code ;
  \item l'application de test m\'ecanique, pr\'esent\'ee \`a la section \ref{sec:testmeca} ;
  \item des extraits de code comment\'es, fournis dans l'annexe \ref{ann:code} ;
  \item la gestion du capteur et de la liaison s\'erie, pr\'esent\'ee \`a la section \ref{sec:capteur} ;
  \item la gestion du bras robotique et du r\'eseau, pr\'esent\'ee \`a la section \ref{sec:robot} ;
  \item les fonctions pertinentes dans des fichiers \texttt{.txt}, fournies dans \path{C:/WORKSPACE/CSM/CSM/livrables/annexes/} ;
  \item un exemple de fichier de test CSV, fourni dans \path{C:/WORKSPACE/CSM/CSM/livrables/annexes/exemple_test_force_position.csv} ;
  \item la composition de l'\'equipe au d\'ebut du PDF, plac\'ee sur la page de garde.
\end{itemize}

Sur le fond, le projet satisfait bien la logique g\'en\'erale demand\'ee : validation du capteur, validation du robot, pilotage manuel en rep\`ere TOOL, acquisition synchronis\'ee et sauvegarde des donn\'ees. En revanche, le point ``asservissement en force'' ne peut pas \^etre pr\'esent\'e comme une boucle PID finalis\'ee, car le code actuel met plut\^ot en place une logique de \emph{dodge} et de r\'eaction \`a l'effort qu'un asservissement acad\'emique complet.

\section{Description rapide de l'IHM}
\label{sec:ihm}
L'application \texttt{RobotForceIntegration} repose sur une IHM WinForms organis\'ee en blocs fonctionnels, et non sous forme d'onglets. Cette structure est importante \`a souligner, car elle refl\`ete l'organisation r\'eelle de l'interface dans \path{C:/WORKSPACE/CSM/CSM/demo/csharp_studio/RobotForceIntegration/Form1.Designer.cs}. Le premier bloc regroupe la connexion robot, l'activation du mouvement, la sensibilit\'e de collision et le contr\^ole de l'autocollision. Un second bloc affiche l'\'etat du robot \`a travers les positions articulaires et cart\'esiennes. Deux autres blocs permettent le jogging manuel : l'un dans le rep\`ere TOOL pour les axes X, Y et Z, l'autre au niveau des articulations.

La partie capteur est elle aussi int\'egr\'ee \`a la m\^eme fen\^etre. L'utilisateur peut y choisir le port s\'erie, les param\`etres de communication, ouvrir ou fermer le port, envoyer une commande libre, ou effectuer une tare. Une zone d'acquisition permet ensuite de choisir le fichier CSV, de d\'emarrer l'essai et d'enregistrer les mesures. Enfin, la zone \emph{Dodge} param\`etre la r\'eaction \`a l'effort, et un journal d'\'ev\'enements retrace les actions importantes de l'application.

Cette IHM est coh\'erente avec un contexte de laboratoire. Elle centralise les op\'erations utiles sans multiplier les fen\^etres, tout en gardant une logique claire entre configuration, contr\^ole manuel, acquisition automatique et diagnostic par le journal.

\section{Gestion du capteur et de la liaison s\'erie}
\label{sec:capteur}
La brique capteur est fond\'ee sur \texttt{System.IO.Ports.SerialPort}. Dans \texttt{RobotForceIntegration}, les param\`etres de communication sont configur\'es dynamiquement par \texttt{ConfigureSerialPort()}, qui applique le port, le baudrate, la parit\'e, les bits de stop, les bits de donn\'ees ainsi que les timeouts de lecture et d'\'ecriture. Le choix par d\'efaut observ\'e dans l'application est compatible avec le besoin exprim\'e pour le capteur : 115200 bauds, format 8N1, et un port COM s\'electionnable depuis l'IHM.

L'acquisition de la force suit ensuite une logique classique mais robuste. L'application envoie une commande au capteur, principalement \texttt{"M"}, puis lit les donn\'ees re\c{c}ues par blocs gr\^ace \`a l'\'ev\'enement \texttt{DataReceived}. Un buffer texte reconstruit les lignes compl\`etes, ce qui \'evite les erreurs de parsing dues \`a des messages s\'erie partiels. La fonction \texttt{TryParseForce()} applique ensuite une expression r\'eguli\`ere pour extraire une valeur num\'erique et la convertir avec une culture invariante.

Le projet \texttt{RS232-PC} joue ici un r\^ole important. Il sert d'application de validation de la cha\^\i ne s\'erie ind\'ependamment du robot. On y retrouve la d\'etection des ports COM, la configuration des param\`etres, l'envoi de commandes, la r\'eception des lignes, et une sauvegarde CSV simple. Cette seconde application constitue donc une brique de test tr\`es utile du point de vue p\'edagogique, car elle permet de qualifier le capteur avant m\^eme d'aborder l'int\'egration robotique.

\section{Gestion du bras robotique et du r\'eseau}
\label{sec:robot}
La communication avec le bras xArm6 est assur\'ee par une architecture en deux niveaux. Le premier niveau est le wrapper \texttt{XArmAPI.cs}, qui d\'eclare les fonctions import\'ees depuis \texttt{xarm.dll}. Le second niveau est la classe \texttt{Robot.cs}, qui encapsule ces appels natifs et maintient un \'etat interne plus s\^ur : connexion, \emph{motion enable}, sensibilit\'e de collision, autocollision, vitesse articulaire et vitesse cart\'esienne.

La connexion se fait \`a partir d'une adresse IP. La m\'ethode \texttt{Create()} cr\'ee ou r\'eactive l'instance robot, s\'electionne la bonne instance active, puis initialise les offsets, la lecture de la pose courante et plusieurs param\`etres de s\'ecurit\'e. Une fois la connexion \'etablie, le robot est pr\'epar\'e au mouvement par l'encha\^\i nement \texttt{EnableMotion(true)} puis \texttt{PrepareMotion()}, cette derni\`ere m\'ethode pla\c{c}ant explicitement le robot en mode position et en \'etat pr\^et.

Le pilotage manuel demand\'e dans le cahier des charges est bien pr\'esent, et m\^eme plus avanc\'e que le minimum attendu. Le besoin initial mentionnait au moins deux boutons \texttt{+X} et \texttt{-X} pour un d\'eplacement de 10 mm en rep\`ere TOOL. Le code actuel propose non seulement ces deux commandes, mais aussi les directions Y et Z, avec un pas r\'eglable depuis l'IHM. La fonction \texttt{MoveToolRelative()} de \texttt{Robot.cs} r\'ealise ce comportement en transformant un incr\'ement \((dx,dy,dz)\) en pose outil relative envoy\'ee \`a l'API du robot.

\section{Couplage Robot/Force}
\label{sec:couplage}
Le couplage entre les deux sous-syst\`emes constitue le c\oe ur du projet. Dans le code, il repose sur \texttt{timerCMD\_Tick}, une boucle d'acquisition p\'eriodique r\'egl\'ee \`a 300 ms. \`A chaque cycle, l'application v\'erifie d'abord que le robot est op\'erationnel et que la liaison s\'erie du capteur est ouverte. Si l'une de ces conditions n'est plus respect\'ee, l'essai est interrompu proprement.

Lorsque les deux briques sont disponibles, la boucle applique un d\'eplacement relatif de l'outil en Z par appel \`a \texttt{xARM.MoveToolRelative(0.0F, 0.0F, GetToolStep(), true)}. La position Z courante est ensuite lue depuis le robot et stock\'ee dans \texttt{pendingPositionZ}. L'application envoie alors la commande \texttt{"M"} au capteur afin de demander une mesure. Lorsqu'une ligne valide est re\c{c}ue et analys\'ee, la force mesur\'ee est associ\'ee \`a la position Z pr\'ec\'edemment m\'emoris\'ee, puis l'ensemble est ajout\'e \`a une liste de mesures.

Ce fonctionnement est s\'equentiel : mouvement, lecture de position, demande de mesure, puis stockage. Il ne s'agit pas encore d'une loi de commande complexe, mais d'un couplage temporel rigoureux entre d\'eplacement et acquisition, qui r\'epond d\'ej\`a au besoin d'un essai instrument\'e.

\section{Asservissement en force}
\label{sec:asservissement}
Le cahier des charges mentionne un asservissement en force, ce qui impose une lecture honn\^ete du code. Dans sa forme actuelle, le projet ne met pas en \oe uvre un PID discret explicite pilotant le mouvement en fonction d'une consigne de force. Aucune structure de type erreur/int\'egrale/d\'eriv\'ee n'appara\^\i t dans la boucle \`a 300 ms, et le d\'eplacement automatique de l'essai reste essentiellement impos\'e.

En revanche, le code int\`egre une logique r\'eelle de r\'eaction \`a l'effort, concentr\'ee dans le mode \emph{Dodge}. Cette logique comporte plusieurs \'el\'ements pr\'eparatoires \`a un futur asservissement : filtrage exponentiel de la force, seuil de d\'etection, calcul d'une vitesse cart\'esienne de retrait, et activation d'un mode de vitesse continue du robot. Le \emph{dodge} n'est donc pas un simple artifice d'IHM. Il constitue une premi\`ere forme de comportement pilot\'e par l'effort, mais orient\'ee vers la protection, l'\'evitement d'obstacle et l'assistance manuelle plut\^ot que vers une r\'egulation acad\'emique de force.

Du point de vue de la conformit\'e au cahier des charges, cette distinction doit \^etre formul\'ee clairement. Le projet fournit une base logicielle compatible avec un futur asservissement en force Z, mais l'impl\'ementation actuelle doit \^etre pr\'esent\'ee comme une logique partielle de r\'eaction \`a l'effort et non comme un PID finalis\'e.

\section{Application de test m\'ecanique}
\label{sec:testmeca}
L'application de test m\'ecanique correspond au mode d'acquisition automatique de \texttt{RobotForceIntegration}. Son principe est simple et coh\'erent avec un premier niveau d'instrumentation : \`a chaque p\'eriode de 300 ms, le robot avance en Z, la force est demand\'ee au capteur, puis le couple \{force, position Z\} est archiv\'e. Ce fonctionnement permet d'obtenir une relation force/d\'eplacement exploitable pour un essai de contact progressif ou de compression.

La sauvegarde des mesures est r\'ealis\'ee en CSV. Le code produit un fichier avec les colonnes \texttt{Timestamp}, \texttt{Force} et \texttt{PositionZ}, ce qui facilite une reprise des donn\'ees dans un tableur ou un outil de traitement externe. Un m\'ecanisme de repli est m\^eme pr\'evu si le fichier cible est verrouill\'e : l'application g\'en\`ere alors automatiquement un nouveau nom de fichier horodat\'e afin d'\'eviter la perte des mesures.

Cette application de test m\'ecanique est donc d\'ej\`a pertinente sur le plan p\'edagogique et exp\'erimental. Elle n'est pas encore une boucle d'asservissement en force compl\`ete, mais elle fournit un environnement robuste pour valider la liaison capteur, le mouvement robotique et la synchronisation des deux.

\section{Bilan critique}
Le code fourni montre une int\'egration s\'erieuse des deux briques mat\'erielles. La partie robot est bien encapsul\'ee, la partie s\'erie est raisonnablement robuste, et l'IHM r\'eunit la plupart des fonctions attendues pour un usage de laboratoire. Le projet satisfait le besoin de d\'eplacements en rep\`ere TOOL, y compris au-del\`a du minimum attendu, et propose un enregistrement CSV directement exploitable.

La principale limite du projet tient \`a la fronti\`ere entre \emph{essai instrument\'e} et \emph{asservissement}. Le code actuel est tr\`es cr\'edible pour la premi\`ere partie, mais il reste incomplet pour la seconde. Il ne faut donc pas survendre l'application : la r\'eaction \`a l'effort existe, mais elle prend la forme d'un \emph{dodge} protecteur et non d'un asservissement PID en force Z.

Une autre limite observable concerne la pr\'esentation visuelle des donn\'ees de force. Le journal d'\'ev\'enements est bien pr\'esent et l'\'etat robot est affich\'e, mais le suivi de force n'est pas encore valoris\'e par une visualisation graphique ou un tableau de mesures dans l'IHM. Une telle extension renforcerait nettement la dimension exp\'erimentale de l'application.

\section{Conclusion}
Le projet CSM Robot + Force constitue une base de travail convaincante pour l'int\'egration d'un capteur exteroceptif et d'un bras robotique dans une application C\# unique. Les objectifs essentiels de la mati\`ere sont atteints sur les volets connexion capteur, connexion robot, mouvement manuel en rep\`ere TOOL, acquisition synchronis\'ee et sauvegarde des essais. Le rendu propos\'e dans ce livrable met en avant ces acquis sans masquer les limites encore pr\'esentes.

Dans une perspective de poursuite du projet, l'\'evolution la plus naturelle consisterait \`a remplacer le d\'eplacement impos\'e en Z par une commande calcul\'ee \`a partir de la force mesur\'ee, avec saturation, fen\^etres de s\'ecurit\'e et, \`a terme, une loi PID discr\`ete. L'architecture actuelle est suffisamment claire pour accueillir cette extension sans remise \`a plat compl\`ete.

\appendix

\section{Extraits de code comment\'es}
\label{ann:code}

Cette annexe rassemble les extraits les plus significatifs du projet. Chaque listing est volontairement court et replac\'e dans sa fonction vis-\`a-vis du cahier des charges.

\subsection{Configuration s\'erie du capteur}
Le bloc suivant montre que la configuration de la liaison s\'erie est bien param\'etrable depuis l'IHM. C'est l'un des points de base exig\'es pour la brique capteur.
\begin{lstlisting}
private void ConfigureSerialPort()
{
    serialPortSensor.PortName = comboBoxSerialPort.Text.Trim();
    serialPortSensor.BaudRate = int.Parse(comboBoxBaudRate.Text, CultureInfo.InvariantCulture);
    serialPortSensor.Parity = (Parity)Enum.Parse(typeof(Parity), comboBoxParity.Text);
    serialPortSensor.StopBits = (StopBits)Enum.Parse(typeof(StopBits), comboBoxStopBits.Text);
    serialPortSensor.DataBits = int.Parse(comboBoxDataBits.Text, CultureInfo.InvariantCulture);
    serialPortSensor.NewLine = "\r\n";
    serialPortSensor.ReadTimeout = 500;
    serialPortSensor.WriteTimeout = 500;
}
\end{lstlisting}

\subsection{Parsing et reconstruction du flux s\'erie}
Cet extrait montre la reconstruction d'une ligne compl\`ete \`a partir d'un buffer, puis l'appel \`a la logique de traitement de la force. Il s'agit d'un point important pour rendre la communication s\'erie fiable.
\begin{lstlisting}
private void ProcessSerialChunk(string chunk)
{
    if (InvokeRequired)
    {
        BeginInvoke(new Action<string>(ProcessSerialChunk), chunk);
        return;
    }

    serialBuffer.Append(chunk);

    while (true)
    {
        string buffer = serialBuffer.ToString();
        int lineBreakIndex = buffer.IndexOf('\n');
        if (lineBreakIndex < 0)
            break;

        string line = buffer.Substring(0, lineBreakIndex).Trim('\r', '\n', '\t', ' ');
        serialBuffer.Remove(0, lineBreakIndex + 1);

        if (line.Length == 0)
            continue;

        HandleSensorLine(line);
    }
}
\end{lstlisting}

\subsection{Extraction de la force et stockage de la mesure}
Ce bloc relie directement la r\'eception du message capteur au stockage du couple \{force, position Z\}. Il concr\'etise le couplage entre la mesure et l'\'etat du robot.
\begin{lstlisting}
private static bool TryParseForce(string line, out float force)
{
    Match match = ReadingRegex.Match(line);
    if (!match.Success)
        match = GenericNumberRegex.Match(line);

    if (!match.Success)
    {
        force = 0.0F;
        return false;
    }

    string numericText = match.Groups["value"].Value.Replace(',', '.');
    return float.TryParse(numericText, NumberStyles.Float, CultureInfo.InvariantCulture, out force);
}

private void HandleSensorLine(string line)
{
    AppendSensorLog("[RX] " + line);

    float force;
    if (!TryParseForce(line, out force))
        return;

    currentForce = force;
    lastForceSampleUtc = DateTime.UtcNow;

    if (!acquisitionRunning || !awaitingForceSample)
        return;

    measures.Add(new Measure
    {
        Timestamp = DateTime.Now,
        Force = currentForce,
        PositionZ = pendingPositionZ
    });
    awaitingForceSample = false;
}
\end{lstlisting}

\subsection{D\'eplacement relatif de l'outil}
La fonction suivante correspond \`a la brique robot minimale demand\'ee pour le rep\`ere TOOL. Elle montre comment un incr\'ement \((dx,dy,dz)\) est traduit en commande robot utilisable par les boutons de jogging.
\begin{lstlisting}
public int MoveToolRelative(float dx, float dy, float dz, bool wait = true, float timeout = NO_TIMEOUT)
{
    float[] pose = { dx, dy, dz, 0.0F, 0.0F, 0.0F };
    return MoveTool(pose, wait, timeout);
}
\end{lstlisting}

\subsection{Boucle d'acquisition Robot/Force \`a 300 ms}
Cet extrait est le plus repr\'esentatif de l'int\'egration globale. Il montre le d\'eplacement en Z, la lecture de la position et l'envoi de la commande \texttt{"M"} dans une m\^eme s\'equence logicielle.
\begin{lstlisting}
private void timerCMD_Tick(object sender, EventArgs e)
{
    if (!IsRobotReady() || !IsSerialReady())
    {
        StopAcquisition(true);
        MessageBox.Show("Robot motion or sensor serial connection is not ready.",
            "Acquisition", MessageBoxButtons.OK, MessageBoxIcon.Warning);
        return;
    }

    if (!ExecuteRobotCommand(() => xARM.MoveToolRelative(0.0F, 0.0F, GetToolStep(), true), "Move Tool +Z"))
    {
        StopAcquisition(true);
        return;
    }

    pendingPositionZ = xARM.GetCurrentPosition()[2];
    awaitingForceSample = true;
    SendSensorCommand("M");
}
\end{lstlisting}

\subsection{Logique de \emph{dodge} pilot\'ee par l'effort}
Cette partie n'est pas un PID, mais elle montre clairement l'existence d'une logique r\'eactive fond\'ee sur la mesure de force. C'est ce point qui justifie de parler d'une base d'asservissement plut\^ot que d'un asservissement complet.
\begin{lstlisting}
private void timerDodge_Tick(object sender, EventArgs e)
{
    if (!IsDodgeRunning())
    {
        SyncDodgeTimerState();
        return;
    }

    if (dodgeAwaitingForceSample &&
        (DateTime.UtcNow - lastDodgeRequestUtc).TotalMilliseconds > DodgeForceSampleTimeoutMilliseconds)
    {
        dodgeAwaitingForceSample = false;
    }

    if (!dodgeAwaitingForceSample)
    {
        dodgeAwaitingForceSample = true;
        lastDodgeRequestUtc = DateTime.UtcNow;
        SendSensorCommand("M");
    }
}
\end{lstlisting}

\subsection{Sauvegarde CSV des mesures}
Le listing suivant r\'epond \`a la demande de sauvegarde en fichier de test. Les trois grandeurs archiv\'ees sont la date, la force et la position Z.
\begin{lstlisting}
private void WriteMeasures(string path)
{
    using (StreamWriter writer = new StreamWriter(path, false, Encoding.UTF8))
    {
        writer.WriteLine("Timestamp,Force,PositionZ");
        foreach (Measure measure in measures)
        {
            writer.WriteLine(
                measure.Timestamp.ToString("O", CultureInfo.InvariantCulture) + "," +
                measure.Force.ToString("G9", CultureInfo.InvariantCulture) + "," +
                measure.PositionZ.ToString("G9", CultureInfo.InvariantCulture));
        }
    }
}
\end{lstlisting}

\section{Annexes fournies avec le livrable}
\label{ann:livrables}
Le dossier \path{C:/WORKSPACE/CSM/CSM/livrables/annexes/} contient les pi\`eces compl\'ementaires demand\'ees pour le ZIP :
\begin{itemize}
  \item \texttt{exemple\_test\_force\_position.csv} : exemple de fichier de test coh\'erent avec le format produit par l'application ;
  \item \texttt{fonctions\_capteur\_serie.txt} : extrait des fonctions les plus pertinentes pour la communication capteur ;
  \item \texttt{fonctions\_robot\_reseau.txt} : extrait des fonctions les plus pertinentes pour la communication robot ;
  \item \texttt{fonctions\_couplage\_robot\_force.txt} : extrait des fonctions qui relient mesure et mouvement dans l'application principale ;
  \item \texttt{README\_livrable.txt} : rappel de la structure propos\'ee pour le ZIP et de la g\'en\'eration automatique du PDF.
\end{itemize}

\end{document}
